\documentclass[12pt,a4paper]{article}

% ── Packages ──
\usepackage[utf8]{inputenc}
\usepackage[T1]{fontenc}
\usepackage{amsmath,amssymb}
\usepackage{graphicx}
\usepackage{booktabs}
\usepackage{tabularx}
\usepackage[bookmarksopen,bookmarksnumbered]{hyperref}
\usepackage{xurl}
\usepackage[margin=2.5cm]{geometry}
\emergencystretch=3em
\usepackage{setspace}
\usepackage[numbers,sort&compress]{natbib}
\usepackage{lineno}
\onehalfspacing
\linenumbers

\renewcommand{\topfraction}{0.95}
\renewcommand{\bottomfraction}{0.95}
\renewcommand{\textfraction}{0.08}
\renewcommand{\floatpagefraction}{0.75}

% ── Title ──
\title{\textbf{World Model for Brain-Blood Mapping and Virtual Knockout Validation Uncovers Drug Repurposing Opportunities and Diagnostic Biomarker Candidates in Alzheimer's Disease}}

\author{Zhongyuan Dong$^{1,\dagger}$, Xuanlin Meng$^{1,\dagger,\ddagger}$ and Lianghua Wang$^{1,2,*}$\\[1em]
\small $^{1}$Department of Biochemistry and Molecular Biology,\\
\small College of Basic Medical Sciences, Naval Medical University,\\
\small Shanghai 200433, China\\[0.5em]
\small $^{2}$Key Laboratory of Biosafety Defense (Naval Medical University),\\
\small Ministry of Education, Shanghai 200433, China\\[0.5em]
\small $^{\dagger}$These authors contributed equally\\[0.3em]
\small $^{\ddagger}$Co-corresponding author: xuanlinmeng@126.com (X.M.)\\[0.3em]
\small $^{*}$Corresponding author: lhwang@smmu.edu.cn (L.W.)}
\date{}

\begin{document}
\maketitle

% ======================================================================
\section*{Abstract}
How molecular signals propagate between the brain and peripheral blood in neurodegeneration is unclear. We construct a \emph{molecular world model} that learns the continuous transformation between brain and blood molecular compartments, using a Neural Ordinary Differential Equation (Neural ODE) as its differentiable engine and Sinkhorn Optimal Transport for distribution matching. The defining capability of this world model is that its learned transformation can be differentiated, yielding Jacobian sensitivity matrices that predict how perturbing any brain-end gene alters blood-end molecular state. Using the 5xFAD mouse Alzheimer's model as the discovery cohort, with 17{,}609 one-to-one orthologs mapped from mouse to human gene symbols, the Jacobian sensitivity network achieves concordance with independent functional knockout validation (AUC\,=\,0.893), far exceeding linear baselines (0.560), and transfers across species: the mouse-trained model produces non-trivial transformations of human ADNI transcriptome data (1,961 of 2,000 model genes matched). Tri-method virtual knockout (GenKI, tissue-specific PPI propagation, and SCENIC gene regulatory network perturbation) nominates directional cross-tissue edges across two omics layers, leading to three highlighted translational outputs. \emph{Diagnostic biomarker}: \textit{MAPT} emerges as a dual-qualified biomarker, diagnostic AUC\,=\,0.719 and survival HR\,=\,1.14 ($p$\,=\,0.003), passing GenKI, PPI propagation, diagnostic, and prognostic validation in the proteomics forward layer. \emph{Drug repurposing}: \textit{CSF3R} passes all four validation tiers in the transcriptomics reverse layer (GenKI + SCENIC edge concordance + survival HR\,=\,1.13, $p$\,=\,0.028 + Phase II drugs filgrastim/pegfilgrastim), while \textit{MMP9} (HR\,=\,1.13, $p$\,=\,0.021; Phase III drug andecaliximab) passes GenKI + survival + clinical drug triage. The two omics layers are complementary: proteomics for diagnostics, transcriptomics for therapeutics.

% ======================================================================
\section{INTRODUCTION}
% ======================================================================

Neurodegenerative diseases are increasingly understood not as isolated pathologies of the brain but as systemic conditions in which central neuropathology and peripheral immune dysfunction are coupled through a bidirectional molecular crosstalk axis \citep{bettcher2021}. In Alzheimer's disease (AD), epidemiological, neuroimaging, and multi-omics evidence converge on a feed-forward loop of neuroinflammation: peripheral blood mononuclear cells from AD patients exhibit pro-inflammatory transcriptional states that correlate with cognitive decline \citep{gate2024, xu2023}, while central amyloid-$\beta$ and tau pathology drive systemic immune responses through blood--brain barrier disruption \citep{sweeney2018, nation2019}, damage-associated molecular patterns \citep{venegas2017}, and innate immune activation \citep{hessel2023}. This cross-tissue interplay is mechanistically central to disease progression, yet the analytical tools used to study it remain predominantly discrete and correlation-based: co-expression and differential-expression analyses operate within single tissues, and cross-tissue relationships are reduced to pairwise correlations that cannot capture the continuous, multivariate transformation of molecular states between compartments \citep{wang2022}.

What is missing is a framework that can learn the continuous distributional transformation from one tissue's molecular space to another's, yielding a principled, model-grounded basis for predicting how perturbations in one compartment propagate to the other. In reinforcement learning, a \emph{world model} learns a predictive simulation of the environment that supports planning and intervention \citep{ha2018}; by direct analogy, a \emph{molecular world model} should learn the transformation rules governing how molecular states propagate across anatomical compartments, enabling in silico perturbation experiments. The defining requirement of such a world model is differentiability: the learned transformation must be differentiable so that its Jacobian can be computed, yielding sensitivity matrices that quantify, for every gene pair, how a unit perturbation in the source tissue alters expression in the target. This Jacobian is the computational analog of a perturbation screen, producing cross-tissue regulatory edges that can be experimentally validated.

We construct such a molecular world model using a Neural Ordinary Differential Equation (Neural ODE) \citep{chen2018} as its differentiable engine, regularized by Sinkhorn Optimal Transport (OT) \citep{cuturi2013} for distribution matching. The Neural ODE parameterizes the derivative of a continuous-time dynamical system with a neural network, learning the nonlinear manifold through which molecular signals propagate across the brain--blood interface without requiring sample-level pairing. Crucially, the ODE formulation makes the entire mapping differentiable: the Jacobian of the integrated trajectory can be computed via automatic differentiation, a capability that no discrete OT or linear regression method provides. This world model is not merely a mapping function but a \emph{simulator} of cross-tissue molecular dynamics, supporting trajectory interpolation, reversibility, and, most importantly, sensitivity analysis for virtual perturbation experiments.

A central design choice is the use of the 5xFAD transgenic mouse model \citep{oakley2006} as the discovery cohort. The 5xFAD model expresses five familial AD mutations and develops amyloid pathology, neuroinflammation, and cognitive deficits by 4--6 months of age \citep{jawhar2012}, providing a biologically relevant system in which paired brain and bone marrow single-cell RNA-seq data can be obtained from the same animals, a pairing that is ethically and practically impossible at scale in human subjects. We map mouse genes to their human one-to-one orthologs using the Mouse Genome Informatics (MGI) homology database \citep{eppig2015}, yielding 17{,}609 orthologous gene pairs, and validate all downstream findings in independent human AD cohorts (ADNI).

Virtual knockout (VK) provides orthogonal computational support for the Jacobian-derived edges and assigns their directionality. We employ two methodologically independent VK approaches: GenKI, a variational graph autoencoder (VGAE) that measures latent-space displacement upon gene deletion \citep{senuma2023}, and, for proteomics, network propagation on a tissue-specific protein--protein interaction (PPI) atlas \citep{lamantrip2025} that captures brain- and blood-specific protein associations; for transcriptomics, SCENIC-based gene regulatory network (GRN) perturbation \citep{aibar2017, moerman2019} infers regulon activity shifts upon gene knockout. Forward VK (knock out brain-end genes, observe blood-end response) nominates diagnostic biomarker candidates; reverse VK (knock out blood-end genes, observe brain-end response) nominates therapeutic target candidates.

We show that the molecular world model learns biologically interpretable cross-tissue mappings; that its Jacobian sensitivity network is the only approach whose predictions are concordant with independent functional knockout validation (AUC\,=\,0.893 vs.\ 0.560 for linear alternatives); that the mouse-trained model transfers to human data without retraining; and that tri-method virtual knockout yields three highlighted translational targets: \textit{MAPT} (dual-qualified diagnostic biomarker, AUC\,=\,0.719, HR\,=\,1.14, $p$\,=\,0.003), \textit{CSF3R} (four-tier validated drug target, filgrastim Phase II), and \textit{MMP9} (Phase III repurposing candidate, andecaliximab). These findings establish the molecular world model as a reusable computational paradigm for cross-tissue drug and biomarker discovery.

% ======================================================================
\section{RESULTS}
% ======================================================================

% ─── Figure 1 ───
\subsection{A Molecular World Model with a Neural ODE Engine Learns Continuous Cross-Tissue Mappings}

We constructed a molecular world model for the brain--blood axis using a Neural ODE engine with Sinkhorn OT regularization (Figure~\ref{fig:framework}; Methods), applied to two molecular layers: transcriptomics (5xFAD mouse brain and bone marrow single-cell RNA-seq, 10{,}000 cells $\times$ 2{,}000 highly variable genes after ortholog mapping) and proteomics (ADNI plasma and CSF NULISA proteomics, 134 targeted proteins). The world model learned a continuous, differentiable transformation between compartments: given a blood-end molecular state $\mathbf{x}_{\text{blood}}$, the ODE engine integrates the learned vector field $f_\theta$ to predict the corresponding brain-end state $\mathbf{z}(1) \approx \mathbf{x}_{\text{brain}}$. Both models converged stably: the Sinkhorn OT loss decreased monotonically over 100 epochs, and terminal distributions visually overlapped with target distributions in PCA-projected spaces.

\textbf{Benchmark proves world-model-exclusive capabilities.}
To demonstrate that the Neural ODE engine provides capabilities unavailable to simpler alternatives, we conducted a benchmark comparing four methods on held-out validation data: (1) Neural ODE (proposed), (2) Direct OT Map (Sinkhorn coupling with barycentric projection), (3) Linear Ridge Regression, and (4) Identity baseline (no mapping) (Table~\ref{tab:benchmark}).

In the \emph{fair end-to-end comparison}, all four methods were evaluated on identical inputs. Neural ODE achieved an OT distance of 3{,}441, a 13.3\% improvement over the identity baseline (3{,}971), while Ridge regression catastrophically failed (OT 6{,}243), confirming that linear transformations cannot capture nonlinear cross-tissue mappings. Direct OT achieved the lowest OT distance (1{,}986), as expected since Sinkhorn directly optimizes this objective, but its gene--gene correlation-structure error was the worst of all methods (0.95 vs.\ 0.21 for identity), confirming distribution-level reordering without preserving multivariate regulatory structure.

In the \emph{world-model-exclusive comparison}, we quantified three capabilities unique to a differentiable continuous model:
\begin{enumerate}
\item \textbf{Jacobian derivability.} The Jacobian sensitivity network's concordance with independent GenKI functional knockout validation yielded AUC\,=\,0.893 for Neural ODE vs.\ 0.560 for Ridge; Direct OT is non-differentiable and cannot produce a Jacobian.
\item \textbf{Continuous trajectory.} Evaluating the ODE at intermediate time points $t \in [0,1]$ revealed a smooth monotonic transition of OT distance to the blood compartment (from 3{,}979 to 3{,}424), with no discontinuities; Direct OT and Ridge provide only two-endpoint maps with no intermediate states.
\item \textbf{Reversibility.} A blood$\rightarrow$brain$\rightarrow$blood round-trip reconstruction achieved Pearson $r$\,=\,0.974 (MSE\,=\,0.084); OT barycentric maps are irreversible.
\end{enumerate}

These results establish that the molecular world model is not merely an incremental mapping improvement but provides qualitatively different capabilities, specifically, a differentiable Jacobian that grounds downstream edge discovery, that no alternative method can offer. The Jacobian AUC of 0.893, confirming that the sensitivity network derived purely from the differentiable vector field is concordant with independent biological perturbation data, is the single most critical benchmark result: it validates the entire downstream pipeline by demonstrating that the world model's learned transformation captures biologically meaningful cross-tissue regulatory structure, not merely statistical distributional similarity.

\textbf{Cross-species transferability.}
To test whether the mouse-trained world model captures evolutionarily conserved regulatory structure, we evaluated it on human ADNI bulk transcriptome data (745 subjects). Of the 2{,}000 model genes, 1{,}961 (98\%) were present in the human microarray. The model produced a non-trivial transformation of human expression data (input--output OT distance\,=\,3{,}838, versus 0 for identity), indicating that it does not simply copy the input but applies a learned brain--blood transformation. The cross-species OT distance between the model's human prediction and real mouse brain expression was 69{,}179, and light fine-tuning (20 epochs) reduced it to 69{,}065 while preserving mouse performance (OT 94{,}122 $\rightarrow$ 93{,}596), indicating that the learned transformation is partially transferable across species without catastrophic forgetting. This cross-species transferability supports the use of the 5xFAD mouse model as a discovery cohort whose findings can be validated in human clinical data.

\begin{table}[htbp]
\centering\small
\caption{\textbf{World-model benchmark: fair end-to-end comparison and exclusive capabilities.}}
\label{tab:benchmark}
\begin{tabular}{lcccc}
\hline
Method & OT distance $\downarrow$ & Corr err.\ $\downarrow$ & Jacobian AUC $\uparrow$ & Round-trip $\uparrow$ \\
\hline
Identity (baseline) & 3{,}971 & 0.21 & --- & --- \\
\textbf{Neural ODE (ours)} & \textbf{3{,}441} & 0.31 & \textbf{0.893} & \textbf{0.974} \\
Direct OT Map & 1{,}986 & 0.95 & N/A & N/A \\
Linear Ridge & 6{,}243 & 0.04 & 0.560 & --- \\
\hline
\end{tabular}
\end{table}

\begin{figure}[htbp]
\centering
\includegraphics[width=\textwidth,height=0.80\textheight,keepaspectratio]{figures/1.pdf}
\caption{\textbf{World-Model Framework and Benchmark.}
\textbf{(a)} Schematic: world-model framework overview (vertical card stack).
\textbf{(b)} 5xFAD single-cell data overview: UMAP of brain and bone marrow cells colored by tissue and genotype.
\textbf{(c)} Mouse-to-human ortholog mapping (Sankey diagram; 17{,}609 one-to-one orthologs).
\textbf{(d)} Neural ODE training convergence: Sinkhorn OT loss over 100 epochs per omics layer.
\textbf{(e)} Fair end-to-end mapping comparison: OT distance by method (Neural ODE vs.\ Direct OT vs.\ Ridge vs.\ Identity).
\textbf{(f)} Jacobian derivability: AUC concordance with functional knockout validation (Neural ODE 0.893 vs.\ Ridge 0.560).
\textbf{(g)} Continuous trajectory: OT distance at intermediate ODE time points $t \in [0,1]$.
\textbf{(h)} Edge confidence metric distributions (stability, SNR, directional consistency) per omics layer from 100 bootstrap replicates.
\textbf{(i)} Transcriptomics top cross-tissue Jacobian sensitivity edges (lollipop).
\textbf{(j)} Proteomics top cross-tissue sensitivity edges (lollipop).
\textbf{(k)} Metabolomics top cross-tissue sensitivity edges (lollipop).}
\label{fig:framework}
\end{figure}

% ─── Figure 2 ───
\subsection{Jacobian Sensitivity Analysis Reveals Cross-Tissue Edges with Multi-Metric Confidence}

From each trained Neural ODE model, we computed the Jacobian sensitivity matrix at the source distribution mean, quantifying how a perturbation of brain molecule $i$ propagates to blood molecule $j$ (Figure~\ref{fig:edges}; Methods). Applying a 95th-percentile threshold on absolute Jacobian values yielded initial sensitivity networks of 198{,}000 cross-tissue edges for transcriptomics and 740 for proteomics, the larger transcriptomic network reflecting the higher dimensionality and cellular heterogeneity of single-cell data.

To quantify edge reliability beyond raw Jacobian magnitude, we computed three complementary confidence metrics from 100 bootstrap replicates: stability (bootstrap agreement on edge presence), signal-to-noise ratio, and directional consistency (forward vs.\ reverse agreement). After confidence and strength filtering, 14{,}244 transcriptomic and 216 proteomic cross-tissue edges were retained.

In the proteomics layer, the strongest cross-tissue edges are directly biologically interpretable. The top-ranked edge connects the tau-phosphorylation pair pTau-231 and pTau-181, followed by MAPT$\leftrightarrow$pTau-181 and amyloid-$\beta$ fragment pairs. Several neuro-peptide axes also emerge, including VEGFD$\leftrightarrow$NPY and BDNF$\leftrightarrow$IL7, linking angiogenic, neurotrophic, and inflammatory signaling across the blood--brain interface. Bidirectional Jacobian computation revealed substantial asymmetry: the median absolute forward Jacobian was approximately 2.8-fold higher than the reverse, suggesting a directional bias that may reflect blood--brain barrier constraints.

\begin{figure}[htbp]
\centering
\includegraphics[width=\textwidth,height=0.80\textheight,keepaspectratio]{figures/2.pdf}
\caption{\textbf{Cross-Tissue Edge Discovery via Jacobian Sensitivity.}
\textbf{(a)} Schematic: Jacobian sensitivity analysis workflow.
\textbf{(b)} Proteomics Jacobian sensitivity heatmap (20$\times$20 hub genes, no labels).
\textbf{(c)} Transcriptomics Jacobian sensitivity heatmap (20$\times$20 hub genes, no labels).
\textbf{(d)} Brain hub centrality ranking per omics layer (3 horizontal bar subplots: Transcriptomics, Proteomics, Metabolomics).
\textbf{(e)} Blood-end eigengene disease-gene ranking per omics layer.
\textbf{(f)} Edge-count reduction through the filtering pipeline (raw $\rightarrow$ confidence-filtered $\rightarrow$ hub-overlap).
\textbf{(g)} Hub--disease gene overlap per omics layer.
\textbf{(h)} Transcriptomics prescreening scatter (mean expression vs.\ variance; selected VK candidate genes highlighted).}
\label{fig:edges}
\end{figure}

% ─── Figure 3 ───
\subsection{Dual-Method Virtual Knockout Validates Cross-Tissue Edges and Assigns Directionality}

To provide orthogonal computational support and assign directionality, we performed virtual knockout (VK) using two methodologically independent approaches (Figure~\ref{fig:vk}; Methods).

\textbf{GenKI (VGAE-based) virtual knockout.}
GenKI constructs a PCNet gene regulatory network from the single-cell data and trains a variational graph autoencoder (VGAE) to learn a latent representation; knockout is simulated by zeroing the target gene's expression, re-training the VGAE, and measuring the Kullback--Leibler divergence between wild-type and knockout latent distributions \citep{senuma2023}. In transcriptomics, 55 of 74 forward-direction brain-end knockouts ($74\%$) and 36 of 45 reverse-direction blood-end knockouts ($80\%$) produced significant downstream effects (Cohen's $d > 0.8$, $p < 0.05$ vs.\ expression-matched negative controls). In proteomics, 8 of 10 forward ($80\%$) and 4 of 25 reverse ($16\%$) GenKI knockouts were significant.

\textbf{Tissue-specific PPI propagation (proteomics).}
For proteomics, we employed a complementary VK based on random walk with restart (RWR) on a tissue-specific protein--protein interaction atlas \citep{lamantrip2025} derived from 7{,}811 proteomic samples across 11 human tissues. Knockout was simulated by removing the target protein from the tissue-specific PPI network and recomputing the RWR steady state. All 9 forward and all 13 reverse proteomics targets showed significant propagation effects (379--679 forward and 434--879 reverse significant downstream targets, $Z > 2$). The union of GenKI and PPI reverse targets comprised 13 therapeutic candidates including APP, IL1B, ICAM1, and TREM1.

\textbf{SCENIC-based GRN perturbation (transcriptomics).}
For transcriptomics, a third independent VK method based on gene regulatory network (GRN) perturbation was applied using the SCENIC framework \citep{aibar2017, moerman2019}. GRNBoost2 \citep{moerman2019} inferred a gene regulatory network from the 5xFAD single-cell expression matrix (3{,}000 cells $\times$ 2{,}078 genes), yielding 136 regulons (transcription factor modules of $\geq$5 co-expressed targets, median size 8). For each knockout gene, AUCell \citep{aibar2017} regulon activity scores were computed before and after zeroing the target gene's expression; the knockout effect was quantified as the mean absolute change in regulon activity across all 136 regulons, with $Z > 1.96$ defining significant perturbation. All 55 forward and all 36 reverse transcriptomics knockouts produced significant downstream effects (mean 8.2 and 7.3 significant regulons per KO, respectively; range 3--13).

Because SCENIC was applied to the same set of KO genes nominated by GenKI, it provides a methodological cross-validation: an independent perturbation principle (GRN regulon activity shift vs.\ VGAE latent-space displacement) applied to the same gene set. To assess edge-level concordance, we examined GenKI-significant edges (KO$\rightarrow$target pairs with $Z > 2$) and tested whether the target's regulon activity was also significantly perturbed in SCENIC upon KO of the same gene. Among 1{,}684 GenKI-significant forward edges, 531 had a corresponding SCENIC regulon measurement, of which 23 (4.3\%) were also significant in SCENIC ($|Z| > 1.96$). In the reverse direction, SCENIC-validated edges converged on \textit{PLA2G7} (lipoprotein-associated phospholipase A2), which was independently nominated as a significant downstream target by six distinct KO genes (MMP8, DACH1, SORL1, CEBPB, CSF3R, IGSF6), representing the strongest cross-method convergence node in the transcriptomics layer. The low overall edge-level concordance (Spearman $\rho$\,=\,0.006 between GenKI and SCENIC effect sizes) reflects the fundamentally different perturbation modalities: GenKI measures global distributional shift in a variational latent space, whereas SCENIC quantifies local regulon activity change in a co-expression network. Both methods confirm that the KO genes produce biologically significant transcriptional perturbation, but they identify different downstream targets, a complementarity that broadens the regulatory coverage of the VK pipeline.

\textbf{Directionality triage.}
Forward VK (brain$\rightarrow$ blood) nominates blood-end genes whose expression changes when a brain-end gene is perturbed, these are diagnostic biomarker candidates, measurable in blood. Reverse VK (blood$\rightarrow$brain) nominates blood-expressed genes whose perturbation impacts brain-end molecular state, these are therapeutic target candidates, amenable to peripheral drug delivery. Table~\ref{tab:validation} summarizes the four-way validation of targets across omics layers and directions.

\begin{table}[htbp]
\centering\small
\caption{\textbf{Multi-layer target validation summary.} Targets passing each validation tier, stratified by omics layer and VK direction.}
\label{tab:validation}
\resizebox{\textwidth}{!}{\begin{tabular}{lllcccl}
\hline
\textbf{Layer} & \textbf{Direction} & \textbf{Target} & \textbf{GenKI} & \textbf{PPI/SCENIC} & \textbf{Clinical} & \textbf{Detail} \\
\hline
\multicolumn{7}{l}{\textit{Fully validated targets (all tiers passed)}} \\
Prot. & Forward & \textbf{MAPT} & \checkmark & PPI \checkmark & AUC\,=\,0.719; HR\,=\,1.14, $p$\,=\,0.003 & Diagnostic + prognostic \\
Trans. & Reverse & \textbf{CSF3R} & \checkmark & SCENIC edge \checkmark & HR\,=\,1.13, $p$\,=\,0.028 & Phase II (filgrastim) \\
\multicolumn{7}{l}{\textit{GenKI + survival + drug (without SCENIC edge concordance)}} \\
Trans. & Reverse & \textbf{MMP9} & \checkmark & SCENIC KO \checkmark & HR\,=\,1.13, $p$\,=\,0.021 & Phase III (andecaliximab) \\
Trans. & Reverse & PGLYRP1 & \checkmark & SCENIC KO \checkmark & HR\,=\,1.19, $p$\,=\,0.001 & Phase 0 \\
Trans. & Reverse & MXD1 & \checkmark & SCENIC KO \checkmark & HR\,=\,1.16, $p$\,=\,0.005 & Phase 0 \\
Trans. & Reverse & S100A6 & \checkmark & SCENIC KO \checkmark & HR\,=\,1.16, $p$\,=\,0.009 & Phase 0 \\
Trans. & Reverse & S100A8 & \checkmark & SCENIC KO \checkmark & HR\,=\,1.16, $p$\,=\,0.013 & Phase 0 \\
Trans. & Reverse & S100A9 & \checkmark & SCENIC KO \checkmark & HR\,=\,1.14, $p$\,=\,0.016 & Phase 0 \\
Trans. & Reverse & TPD52 & \checkmark & SCENIC KO \checkmark & HR\,=\,0.87, $p$\,=\,0.006 & Phase 0 (protective) \\
\multicolumn{7}{l}{\textit{Proteomics reverse (GenKI $\cap$ PPI; no survival significance)}} \\
Prot. & Reverse & REST & \checkmark & PPI \checkmark & AUC\,=\,0.617 & Diagnostic only \\
Prot. & Reverse & VCAM1 & \checkmark & PPI \checkmark & --- & Phase 0 \\
Prot. & Reverse & IGFBP7 & \checkmark & PPI \checkmark & --- & Phase 0 \\
Prot. & Reverse & POSTN & \checkmark & PPI \checkmark & --- & Phase 0 \\
\hline
\end{tabular}}
\end{table}

\begin{figure}[htbp]
\centering
\includegraphics[width=\textwidth,height=0.80\textheight,keepaspectratio]{figures/3.pdf}
\caption{\textbf{Tri-Method Virtual Knockout Validation.}
\textbf{(a)} Schematic: GenKI VGAE, PPI propagation, and SCENIC GRN perturbation workflow.
\textbf{(b)} PPI propagation circular network, forward direction (brain$\rightarrow$blood; KO nodes left, top-5 targets right; edge width = effect score).
\textbf{(c)} PPI propagation circular network, reverse direction (blood$\rightarrow$brain).
\textbf{(d)} GenKI per-KO KL divergence radial bar chart (4 sectors of 90\textdegree{} each: Proteomics-Forward, Proteomics-Reverse, Transcriptomics-Forward, Transcriptomics-Reverse; sector identity in legend).
\textbf{(e)} GenKI confidence scatter (2$\times$2 grid: omics $\times$ direction; $x$ = log expression, $y$ = Cohen's $d$, marker size = $-\log_{10}(p)$, alpha = significant).
\textbf{(f)} SCENIC regulon perturbation per KO (1$\times$2 bar: forward + reverse, top-8 each).
\textbf{(g)} Cross-method concordance (1$\times$2 scatter: GenKI vs.\ PPI for proteomics; GenKI vs.\ SCENIC for transcriptomics; Spearman $\rho$ annotated).
\textbf{(h)} Per-KO target KL divergence distribution (ridge plot, transcriptomics only; KO target vs.\ negative control).
\textbf{(i)} Pathway enrichment of VK-perturbed targets (gene$\rightarrow$pathway Sankey with aligned bubble panel; forward neuronal-projection terms vs.\ reverse immune/trafficking terms).
\textbf{(j)} KO$\times$KO chord diagram of shared perturbation targets (reverse direction; edge weight = number of shared targets).}
\label{fig:vk}
\end{figure}

% ─── Figure 3 supplement: pathway enrichment (kept inside VK section as a paragraph) ───
\textbf{Pathway enrichment of VK perturbation targets.}
To translate the perturbation sets into biological-process terms, we performed Gene Ontology Biological Process over-representation analysis (Enrichr) on the union of significantly perturbed targets ($Z > 2$ across all KOs) per method and direction (Figure~\ref{fig:vk}i; Methods). The enrichment mirrored the diagnostic-vs-therapeutic split: forward (brain$\rightarrow$blood) PPI perturbations (3{,}006 genes) enriched axonogenesis (GO:0007409, adjusted $p = 9.5 \times 10^{-14}$) and neuron-projection guidance, whereas reverse (blood$\rightarrow$brain) perturbations (4{,}678 genes) enriched vesicle-mediated transport (GO:0016192, $p = 2.1 \times 10^{-9}$), ubiquitin-dependent protein catabolism, and, in the GenKI transcriptomics reverse set, regulation of monocyte chemotaxis (GO:0090025, $p = 2.7 \times 10^{-4}$), consistent with the immune-axis therapeutic targets emerging from this layer (CSF3R, MMP9).

% ─── Figure 4 ───
\subsection{Clinical Translation I: Proteomics Virtual Knockout Yields Diagnostic Biomarkers Predicting Cognitive Decline}

Forward-direction VK nominates blood-end proteins whose levels respond to brain-end perturbation, candidate diagnostic biomarkers measurable in plasma. We evaluated these in the ADNI NULISA plasma proteomics cohort ($n = 305$ subjects with baseline CN or AD diagnosis).

\textbf{Network-guided diagnostic panel.}
Among the VK-nominated proteomics targets, \textit{MAPT} achieved the highest single-gene diagnostic AUC of 0.719 for discriminating AD from cognitively normal controls, followed by the merged 4-protein panel (FABP3, IGFBP7, MAPT, NRGN; AUC\,=\,0.696) and CST3 (AUC\,=\,0.630) (Figure~\ref{fig:diagnostic}c, d). The transcriptomics L1-regularized logistic regression panel (LASSO; 17 auto-selected features from 100 VK-nominated genes; Figure~\ref{fig:diagnostic}g) achieved AUC\,=\,0.653 $\pm$ 0.085 (5-fold CV, 5 repeats), while a fixed pre-defined gene panel achieved AUC\,=\,0.629 (Figure~\ref{fig:diagnostic}d). This gap is biologically expected: plasma proteins directly reflect neuronal leakage and central pathology, whereas peripheral blood cell mRNA captures systemic immune states that are indirectly coupled to brain pathology.

\textbf{Prognostic value and aging-time dynamics of diagnostic biomarkers.}
Beyond cross-sectional discrimination, we assessed whether these biomarkers predict \emph{longitudinal} cognitive trajectory using Cox proportional hazards regression in the ADNI cohort ($n$\,=\,620; event = MMSE decline $\geq$ 3 points; APOE4-adjusted). \textit{MAPT} (tau protein) showed the strongest prognostic effect (HR\,=\,1.14 per SD increase, 95\% CI [1.05, 1.24], $p$\,=\,0.003; log-rank $p$\,=\,0.0004), confirmed by Kaplan--Meier analysis splitting subjects at the median MAPT NPQ (Figure~\ref{fig:diagnostic}h; log-rank $p$\,<\,0.001). \textit{TREM1} (myeloid innate immune receptor; HR\,=\,1.11, $p$\,=\,0.011) and \textit{ENO2} (neuron-specific enolase; HR\,=\,0.89, $p$\,=\,0.012) also predicted decline. \textit{MAPT} thus passes both diagnostic (AUC\,=\,0.719) and prognostic (Cox + KM) validation, a dual-qualified biomarker. To locate \emph{when in the lifespan} these targets begin to dysregulate, we projected their expression onto a continuous aging pseudotime computed from Tabula Muris Senis brain single-cell RNA-seq (5{,}000 cells across 3m/18m/24m; diffusion map + DPT rooted at the 3-month centroid; Spearman(age, pseudotime) $\rho = 0.28$, $p = 6 \times 10^{-91}$; Figure~\ref{fig:diagnostic}j; Methods). The 26 ortholog-mapped targets exhibit non-uniform aging dynamics: \textit{MAPT} declines with aging (consistent with cumulative neuronal loss), whereas the therapeutic target \textit{MMP9} rises in parallel with blood--brain-barrier weakening, and several targets (e.g.\ \textit{TGFBR1}, \textit{PYGL}) shift most sharply across the mid-life 18m transition. This dynamic view complements the Cox evidence by showing that the prognostic biomarkers' effects accumulate over the aging trajectory rather than arising abruptly at diagnosis.

\textbf{Transcriptomics diagnostic panel.}
Applying the same pipeline to transcriptomics targets (ADNI gene expression microarray, 305 subjects), the L1-regularized logistic regression panel (LASSO; 17 auto-selected features from 100 VK-nominated genes; Figure~\ref{fig:diagnostic}g) achieved AUC\,=\,0.653 $\pm$ 0.085, and a fixed pre-defined gene panel achieved AUC\,=\,0.629 (Figure~\ref{fig:diagnostic}d), moderate but lower than proteomics. This gap is biologically expected: plasma proteins directly reflect neuronal leakage and central pathology, whereas peripheral blood cell mRNA captures systemic immune states that are indirectly coupled to brain pathology.

\begin{figure}[htbp]
\centering
\includegraphics[width=\textwidth,height=0.80\textheight,keepaspectratio]{figures/4.pdf}
\caption{\textbf{Diagnostic Translation.}
\textbf{(a)} Schematic: diagnostic validation workflow (VK pool $\rightarrow$ ADNI cohort $\rightarrow$ AUC + Cox $\rightarrow$ MAPT validated).
\textbf{(b)} ADNI cohort composition (CN / MCI / AD).
\textbf{(c)} Top single-gene diagnostic AUC (proteomics + transcriptomics).
\textbf{(d)} Panel-level diagnostic AUC comparison: proteomics merged 4-protein panel (FABP3, IGFBP7, MAPT, NRGN) vs.\ transcriptomics L1-regularized logistic regression panel (LASSO, 17 auto-selected features from 100 VK-nominated genes) vs.\ transcriptomics fixed gene set (pre-defined gene panel without feature selection).
\textbf{(e)} ROC curves for proteomics and transcriptomics diagnostic panels.
\textbf{(f)} Proteomics Cox forest plot: 17 genes HR with 95\% CI.
\textbf{(g)} Transcriptomics L1-selected feature coefficients (diverging lollipop; risk vermilion, protective blue).
\textbf{(h)} MAPT Kaplan--Meier curve: high vs.\ low plasma MAPT (ADNI, $n$\,=\,403, log-rank $p$\,<\,0.001).
\textbf{(i)} Transcriptomics Cox forest plot: top-15 genes HR (risk + protective balanced).
\textbf{(j)} Target expression dynamics along an aging pseudotime: 3D surface of 26 validated-target orthologs across the Tabula Muris Senis brain aging trajectory (3m$\rightarrow$18m$\rightarrow$24m); $z$-axis = mean $z$-expression per pseudotime bin.}
\label{fig:diagnostic}
\end{figure}

% ─── Figure 5 ───
\subsection{Clinical Translation II: Transcriptomics Virtual Knockout Identifies Repurposable Drug Targets with Clinical Evidence}

Reverse-direction VK nominates blood-expressed genes whose perturbation impacts the brain-end molecular state, candidate therapeutic targets amenable to peripheral intervention. We evaluated these in two dimensions: (1) whether their expression levels predict disease progression (therapeutic response indicator), and (2) whether existing clinical-stage drugs target them (repurposing potential).

\textbf{Transcriptomics therapeutic targets show prognostic value.}
Among 45 reverse-direction blood-end knockout genes in the 5xFAD transcriptomics layer (36 GenKI-significant), 33 were measurable in the ADNI gene expression microarray. Cox regression (APOE4-adjusted, event = MMSE decline $\geq$ 3) identified 9 genes whose expression levels significantly predict cognitive decline ($p < 0.05$). The top prognostic targets included \textit{PGLYRP1} (HR\,=\,1.19, $p$\,=\,0.001), \textit{IGSF6} (HR\,=\,1.18, $p$\,=\,0.002), \textit{MMP9} (HR\,=\,1.13, $p$\,=\,0.021), \textit{CSF3R} (HR\,=\,1.13, $p$\,=\,0.028), and the S100 alarmin family (\textit{S100A6}, \textit{S100A8}, \textit{S100A9}; all $p < 0.02$). Here, survival significance indicates that the target's expression level influences disease trajectory, a prerequisite for therapeutic modulation.

\textbf{Drug repurposing: three targets pass GenKI + survival + clinical drug triage.}
Multi-database drug mining (OpenTargets, ChEMBL, DGIdb, PubChem) revealed that among the transcriptomics therapeutic targets passing GenKI + survival validation ($p < 0.05$), three have clinical-stage drugs (Phase II or above):
\begin{itemize}
\item \textbf{\textit{MMP9} (matrix metalloproteinase-9):} Phase III inhibitors including andecaliximab (anti-MMP9 antibody), marimastat (broad-spectrum MMP inhibitor), and AZD-1236 (selective MMP-9/12 inhibitor), originally developed for cancer and fibrosis. HR\,=\,1.13 ($p$\,=\,0.021). MMP9 passes GenKI and SCENIC KO validation but not SCENIC edge concordance; its effect is confirmed by both VK methods at the gene level.
\item \textbf{\textit{CXCR2} (C-X-C chemokine receptor 2):} Phase III small-molecule antagonists including danirixin, elubrixin, reparixin, and SX-682, originally developed for COPD, asthma, and cancer immunotherapy. HR\,=\,1.11 ($p$\,=\,0.054). CXCR2 is a key neutrophil chemokine receptor that mediates neuroinflammatory cell recruitment to the brain; its pharmacological antagonism reduces amyloid pathology and cognitive decline in AD mouse models \citep{bakshi2011, ryu2015}.
\item \textbf{\textit{CSF3R} (G-CSF receptor):} Phase II modulators including filgrastim, pegfilgrastim, and lipegfilgrastim (recombinant G-CSF), originally developed for chemotherapy-induced neutropenia. HR\,=\,1.13 ($p$\,=\,0.028). CSF3R passes all four validation tiers: GenKI, SCENIC edge concordance, survival, and clinical drug triage, making it the most comprehensively validated target in this study.
\end{itemize}
The remaining 6 survival-significant targets (PGLYRP1, MXD1, S100A6, S100A8, S100A9, TPD52; all Phase 0 tractability) represent earlier-stage drug discovery opportunities.

\textbf{Proteomics therapeutic targets: intersection vs.\ union strategy.}
Under the GenKI-only intersection strategy (4 targets: IGFBP7, POSTN, REST, VCAM1), none have clinical-stage drugs, these are structural/adhesion proteins without established inhibitor programs. However, expanding to the union of GenKI and PPI propagation VK (13 targets) adds targets with extensive clinical drug pipelines:

\begin{itemize}
\item \textit{TREM1} (triggering receptor expressed on myeloid cells 1): survival-significant (HR\,=\,1.11, $p$\,=\,0.011), Phase II antibody (PY314). TREM1 is a central amplifier of innate immune responses in AD; its inhibition reduces neuroinflammation and amyloid pathology in mouse models \citep{trem1review2024}.
\item \textit{MAPT} (microtubule-associated protein tau): survival-significant (HR\,=\,1.14, $p$\,=\,0.003), Phase II tau-aggregation inhibitors (methylthioninium, LMTX). MAPT is the core pathological protein in AD, and its plasma levels serve as both diagnostic and prognostic biomarker.
\item \textit{APP} (amyloid precursor protein): Phase III anti-amyloid antibodies (Lecanemab, Donanemab, Aducanumab), three recently FDA-approved AD therapies, but did not reach survival significance ($p$\,=\,0.15), consistent with the modest clinical efficacy of anti-amyloid immunotherapy.
\item \textit{IL1B} (interleukin-1$\beta$): Phase III biologics (Canakinumab, Anakinra), anti-inflammatory antibodies originally for autoinflammatory diseases, but did not reach survival significance ($p$\,=\,0.56).
\item \textit{ICAM1} (intercellular adhesion molecule 1): Phase III antibodies (Natalizumab) and antisense oligonucleotides (Alicaforsen), but did not reach survival significance.
\end{itemize}

The GenKI-only intersection strategy thus yields a conservative but pharmacologically empty set, while the union strategy captures clinically validated targets whose inclusion is supported by the orthogonal PPI propagation method. Under the union, \textit{TREM1} and \textit{MAPT} achieve full triage qualification (VK + survival + clinical drug).

\textbf{Complementary omics, complementary translational pathways.}
The proteomics and transcriptomics layers nominate targets in distinct biological pathways: proteomics preferentially captures the tau pathology axis (\textit{MAPT}, \textit{TREM1}), while transcriptomics preferentially captures the immune/inflammation axis (\textit{MMP9}, \textit{S100A8/9}, \textit{CXCR2}). This complementarity reflects the different molecular information captured by each platform: plasma proteins are enriched for neuronal leakage products (central pathology readout), whereas bone marrow single-cell RNA captures active peripheral immune cell states (systemic inflammation readout).

\begin{figure}[htbp]
\centering
\includegraphics[width=\textwidth,height=0.80\textheight,keepaspectratio]{figures/5.pdf}
\caption{\textbf{Therapeutic Translation.}
\textbf{(a)} Schematic: drug repurposing workflow (reverse VK $\rightarrow$ survival filter $\rightarrow$ drug mining $\rightarrow$ tiered ranking).
\textbf{(b)} Transcriptomics reverse-KO Cox forest plot (top-8 risk + top-8 protective genes, balanced; drug phase annotated).
\textbf{(c)} Proteomics Cox forest plot (genes HR with 95\% CI).
\textbf{(d)} Druggable target summary (horizontal bar; top-12 by drug evidence score; colored by clinical phase: Phase 3 vermilion, Phase 2 orange, Tractable blue).
\textbf{(e)} Kaplan--Meier curves for the two Phase III targets: \textit{MMP9} (left) and \textit{CXCR2} (right), high vs.\ low expression (ADNI transcriptomics).
\textbf{(f)} Druggability tier distribution (stacked bar: Proteomics vs.\ Transcriptomics $\times$ tier).
\textbf{(g)} Pathway complementarity (diverging bar: proteomics tau axis vs.\ transcriptomics immune axis).
\textbf{(h)} Drug-evidence bubble chart ($x$ = hazard ratio, $y$ = $-\log_{10}$(Cox $p$), bubble size = drug evidence score, color = clinical phase).}
\label{fig:therapeutic}
\end{figure}

% ======================================================================
\section{DISCUSSION}
% ======================================================================

We have presented a world-model-inspired framework that uses a Neural ODE engine to learn continuous cross-tissue molecular mappings, demonstrated its exclusive analytical capabilities through rigorous benchmarking, and translated its predictions into both diagnostic biomarkers and repurposable drug targets for Alzheimer's disease. Three aspects merit discussion.

\textbf{The molecular world model provides qualitatively different capabilities, not incremental improvements.}
The benchmark analysis reveals that the Neural ODE engine is not simply a better mapping function; it is the \emph{only} approach that yields a differentiable Jacobian sensitivity matrix. This is the critical distinction: the Jacobian transforms a distribution-matching exercise into a predictive perturbation model, enabling in silico knockout experiments whose results are concordant with independent functional validation (AUC\,=\,0.893). Direct OT achieves lower distributional distance but cannot produce a Jacobian (non-differentiable); Ridge regression is differentiable but its linear Jacobian fails to capture nonlinear regulatory structure (AUC\,=\,0.560). The continuous trajectory and reversibility properties further distinguish the world model as a genuine molecular \emph{simulator}, not merely a cross-tissue regressor. Moreover, the mouse-trained world model transfers across species: it produces non-trivial transformations of human ADNI bulk transcriptome data (1,961 of 2,000 model genes matched, 745 subjects), and light fine-tuning reduces the cross-species OT distance while preserving mouse performance, indicating that the learned brain--blood transformation captures evolutionarily conserved regulatory structure rather than species-specific artifacts.

\textbf{Multi-omics complementarity: proteomics for diagnostics, transcriptomics for therapeutics.}
A striking finding is that the two omics layers, though analyzing the same 5xFAD discovery cohort and the same ADNI validation population, nominate targets in different biological pathways with different translational strengths. The proteomics layer excels in diagnostics (\textit{MAPT} single-gene AUC\,=\,0.719; merged panel AUC\,=\,0.696) because plasma proteins are enriched for neuronal origin products (\textit{MAPT}, \textit{ENO2}, \textit{CST3}) that directly reflect central pathology. The transcriptomics layer excels in therapeutics because bone marrow single-cell RNA captures active peripheral immune cell states, including matrix remodeling (\textit{MMP9}), innate immunity (\textit{S100A8/9}), and granulopoiesis (\textit{CSF3R}), whose drug-targeting is more tractable (small-molecule inhibitors vs.\ antibodies for neuronal proteins). This complementarity is not an artifact of data quality but reflects the fundamental biology: different molecular layers interrogate different aspects of the brain--blood axis, and their integration in a single framework captures a broader therapeutic landscape than either alone.

\textbf{Three highlighted translational targets.}
Our framework nominates three targets that warrant prioritized discussion.

\textit{MAPT} (microtubule-associated protein tau) emerges from the proteomics forward layer as the single most validated diagnostic biomarker in this study. It passes GenKI virtual knockout, PPI propagation, diagnostic discrimination (single-gene AUC\,=\,0.719, the highest among all VK-nominated proteomics targets), and Cox prognostic validation (HR\,=\,1.14, $p$\,=\,0.003; Kaplan--Meier log-rank $p$\,<\,0.001). This dual qualification, cross-sectional discrimination and longitudinal prognostic value, is consistent with the explosive growth of plasma p-tau species as the leading blood-based AD biomarkers. Recent studies demonstrate that plasma p-tau217 achieves diagnostic AUC $\sim$0.90 for amyloid pathology \citep{hansson2024}, and novel tau fragments such as MTBR-tau243 identify tau tangle pathology at the MCI stage with strong PET correlation ($\beta$\,=\,0.72, $R^2$\,=\,0.56) \citep{barthelemy2025}. Our framework's identification of \textit{MAPT} from an unbiased cross-tissue VK pipeline, without prior knowledge of the tau biomarker field, provides orthogonal computational support for its clinical priority. The \textit{MAPT} plasma level is measurable on the NULISA proteomics platform used in ADNI, and its prognostic HR of 1.14 per SD increase aligns with published effect sizes for plasma total tau in AD progression cohorts.

\textit{CSF3R} (colony-stimulating factor 3 receptor, G-CSF receptor) is the most comprehensively validated drug target in this study, passing all four tiers: GenKI knockout, SCENIC edge-level concordance, ADNI survival (HR\,=\,1.128, $p$\,=\,0.028), and clinical drug triage (Phase II: filgrastim, pegfilgrastim, lipegfilgrastim). CSF3R is the primary receptor governing granulopoiesis and neutrophil homeostasis. The causal link between CSF3R signaling and AD pathology is supported by multiple lines of evidence: (1) G-CSF treatment decreases brain amyloid burden and reverses cognitive impairment in AD mouse models by attenuating pro-inflammatory cytokine release (TNF-$\alpha$, IL-1$\beta$, IL-12) \citep{tsai2007}; (2) neutrophils, the effector cells of CSF3R signaling, promote AD-like pathology through LFA-1 integrin-mediated CNS infiltration and neutrophil extracellular trap formation \citep{zenaro2015}; (3) a Phase I clinical trial of filgrastim in AD patients (NCT01617577) demonstrated safety and trended toward cognitive improvement. Our framework's identification of CSF3R through the transcriptomics reverse VK layer, which captures bone marrow immune cell states, provides a mechanistically coherent nomination: the same biological axis (granulopoiesis $\rightarrow$ neutrophil infiltration $\rightarrow$ neuroinflammation) is independently captured by the world model's learned brain--blood mapping. The SCENIC edge-level concordance (CSF3R KO significantly perturbs the PLA2G7 regulon, a vascular inflammation marker) further corroborates this target's network-level impact.

\textit{MMP9} (matrix metalloproteinase-9) is the Phase III drug repurposing candidate identified through the transcriptomics reverse layer. It passes GenKI knockout, SCENIC gene-level KO validation, ADNI survival (HR\,=\,1.129, $p$\,=\,0.021), and clinical drug triage (Phase III: andecaliximab, a selective anti-MMP9 monoclonal antibody originally developed for cancer and fibrosis). MMP9 degrades the blood--brain barrier basement membrane and tight junction proteins; plasma MMP9 levels are elevated in AD patients and correlate with BBB permeability \citep{lorenzl2003, viggars2021}. Pharmacological MMP9 modulation improves specific neurobehavioral deficits in AD mouse models \citep{mmp9mouse2021}, and MMP9 has been proposed as a therapeutic target of neuroinflammation in AD \citep{li2025mmp}. While MMP9 does not pass SCENIC edge-level concordance (reflecting the different perturbation modalities of GRN vs.\ VGAE, as discussed above), its validation across GenKI + survival + Phase III drug evidence constitutes strong multi-modal support. Notably, MMP9 and CSF3R are mechanistically linked: neutrophil-driven MMP9 release at the BBB is a downstream effector of CSF3R-mediated granulopoiesis, suggesting these two targets act on the same neuroinflammatory axis at different nodes.

\textbf{The dual meaning of survival analysis.}
A key methodological point is that Cox regression serves different purposes in the diagnostic and therapeutic contexts. For forward-direction biomarker candidates (diagnostic), survival significance indicates \emph{prognostic value}: the biomarker predicts the trajectory of cognitive decline, establishing its clinical utility as a staging or monitoring tool. For reverse-direction therapeutic targets, survival significance indicates \emph{therapeutic response prediction}: the target's endogenous expression level causally influences disease progression, providing human genetic evidence that pharmacological modulation is likely to alter disease trajectory. \textit{MAPT} exemplifies the former (a prognostic biomarker), while \textit{CSF3R} and \textit{MMP9} exemplify the latter (therapeutic targets whose expression levels predict and potentially drive cognitive decline).

\textbf{Limitations.}
The 5xFAD mouse model primarily recapitulates amyloid pathology; tau pathology, neurodegeneration, and human-specific immune features are incompletely represented. Bone marrow is used as a blood surrogate, its cellular composition differs from peripheral blood. The ortholog mapping strategy assumes functional conservation, which holds for most genes but not all. The ADNI microarray platform measures fewer genes than RNA-seq, limiting transcriptomic coverage. Finally, all VK predictions are computational and require experimental validation; the clinical survival and drug evidence provides human corroboration but does not constitute proof of therapeutic efficacy.

\textbf{Future directions.}
The framework is designed as a reusable engine: new cross-tissue datasets can be processed through the same pipeline without modification. Future work will (1) validate MMP9 inhibition in 5xFAD behavioral experiments, (2) extend the framework to incorporate multi-modal data integration (protein post-translational modifications, metabolomics), and (3) build a web interface enabling other researchers to upload cross-tissue data and run the full pipeline, supporting the vision of AI-driven virtual cell discovery \citep{bunne2024}.

% ======================================================================
\section*{Methods}
% ======================================================================

\subsection*{Data availability}
The 5xFAD single-cell RNA-seq dataset is available at GEO (GSE329430). The Tabula Muris Senis dataset is available at figshare (doi:10.6084/m9.figshare.8278114) and the CZ Single-Cell Data Portal. ADNI data (NULISA proteomics, gene expression microarray, DXSUM diagnosis, MMSE, APOE genotyping) are available through the ADNI portal (\url{https://adni.loni.usc.edu}). The tissue-specific PPI atlas is available at BioStudies (S-BSST1423). Drug evidence data were retrieved from OpenTargets (\url{https://platform.opentargets.org}), DGIdb (\url{https://www.dgidb.org}), ChEMBL (\url{https://www.ebi.ac.uk/chembl}), and PubChem (\url{https://pubchem.ncbi.nlm.nih.gov}). Pathway enrichment used the Enrichr GO Biological Process 2023 gene-set library via gseapy. ADNI data are available to qualified investigators upon application through the ADNI Data Portal and are subject to the ADNI Data Use Agreement; they cannot be redistributed by the authors.

\subsection*{Code availability}
All code supporting the findings of this study is available in a public GitHub repository (\url{https://github.com/dong-zhongyuan/AD-Multi-Omics-Integration}) under the MIT license. The repository contains the full analysis pipeline (Neural ODE training, Sinkhorn optimal transport, Jacobian sensitivity analysis, tri-method virtual knockout, diagnostic and survival modelling, cross-species transfer evaluation, aging pseudotime, and pathway enrichment), together with scripts reproducing every main-text figure. A versioned, DOI-minted archive of the exact code used to produce the reported results is deposited at Zenodo (\url{https://doi.org/10.5281/zenodo.XXXXXXX}). The pipeline is implemented in Python~3.10 with PyTorch~2.2, torchdiffeq~0.2.3, geomloss (Sinkhorn), scanpy, and lifelines; pathway enrichment uses gseapy; R-based figures use circlize, ggridges, and ggsankeyfier. All third-party dependencies are open-source and listed in \texttt{environment/requirements.txt} with pinned versions.

% ── Key Resources Table ──
\begin{table}[htbp]
\centering\small
\caption{\textbf{Key Resources.}}
\label{tab:resources}
\resizebox{\textwidth}{!}{\begin{tabular}{lll}
\hline
\textbf{Resource} & \textbf{Source} & \textbf{Identifier} \\
\hline
\multicolumn{3}{l}{\textit{Datasets}} \\
5xFAD brain + bone marrow scRNA-seq & GEO & GSE329430 \\
Tabula Muris Senis brain scRNA-seq & Tabula Muris Senis Consortium & \citep{schaum2019} \\
ADNI NULISA plasma proteomics & ADNI & BSHRI\_PLA\_CSF\_NULISA \\
ADNI gene expression microarray & ADNI & ADNI\_Gene\_Expression\_Profile \\
ADNI longitudinal diagnosis (DXSUM) & ADNI & DXSUM \\
ADNI MMSE scores & ADNI & MMSE \\
ADNI APOE genotyping & ADNI & APOERES \\
Tissue-specific PPI atlas & BioStudies & S-BSST1423 \\
MGI orthology database & MGI & HOM\_MouseHumanSequence.rpt \\
\multicolumn{3}{l}{\textit{Software and Algorithms}} \\
Neural ODE + OT framework & This paper & [GitHub URL] \\
GenKI VGAE virtual knockout & Yang et al.\ 2023 & \citep{senuma2023} \\
Tissue-specific PPI propagation & This paper & [GitHub URL] \\
pySCENIC / GRNBoost2 & Aibar et al.\ 2017 & \citep{aibar2017, moerman2019} \\
MultiXrank & Baptista et al.\ 2022 & \citep{baptista2022} \\
lifelines (Cox regression) & Cam Davidson-Pilon & \url{https://lifelines.readthedocs.io} \\
OpenTargets API & Open Targets & \url{https://platform.opentargets.org} \\
DGIdb (drug-gene interactions) & DGIdb & \url{https://www.dgidb.org} \\
ChEMBL (bioactivity data) & EMBL-EBI & \url{https://www.ebi.ac.uk/chembl} \\
PubChem (bioactivity screening) & NCBI & \url{https://pubchem.ncbi.nlm.nih.gov} \\
Enrichr GO-BP 2023 gene sets & Ma'ayan Lab & \citep{kuleshov2016} \\
scanpy (single-cell analysis, pseudotime) & Wolf et al.\ 2018 & \citep{wolf2018} \\
gseapy / Enrichr (pathway ORA) & Kuleshov et al.\ 2016 & \citep{kuleshov2016} \\
ggsankeyfier (Sankey diagrams) & S'angiamo & \url{https://github.com/SAngiamo/ggsankeyfier} \\
\hline
\end{tabular}}
\end{table}

\subsection*{Experimental Model and Subject Details}

\textbf{5xFAD mice.} Single-cell RNA-seq data were obtained from GEO GSE329430, comprising brain cortex and bone marrow cells from 5xFAD transgenic mice and wild-type littermate controls. Cells were filtered for quality (minimum 200 genes per cell, maximum 20\% mitochondrial reads). Mouse gene symbols were mapped to human one-to-one orthologs using the MGI HOM\_MouseHumanSequence.rpt database, yielding 17{,}609 orthologous gene pairs. After filtering, 10{,}000 cells (5{,}000 brain, 5{,}000 bone marrow) $\times$ 11{,}982 orthologous genes were retained. A separate gene set of 2{,}000 highly variable genes was selected for Neural ODE training; the full gene set (11{,}982 genes) was retained for virtual knockout experiments.

\textbf{ADNI cohort.} The Alzheimer's Disease Neuroimaging Initiative (ADNI) provided plasma proteomics (NULISA platform, 134 proteins), gene expression microarray (745 subjects, 49{,}386 probes), longitudinal diagnosis (DXSUM), MMSE scores, and APOE genotyping. Baseline cognitively normal (CN, diagnosis\,=\,1) and AD (diagnosis\,=\,3) subjects were used for diagnostic analysis ($n$\,=\,305 with both proteomics and expression data). For survival analysis, all subjects with baseline MMSE were included ($n$\,=\,1{,}386).

\subsection*{Method Details}

\subsubsection{Data Preprocessing}

\textbf{5xFAD single-cell preprocessing.} Raw count matrices were loaded and quality-filtered. Highly variable genes were selected using Scanpy's \texttt{seurat} flavor. For Neural ODE training, 2{,}000 HVGs common to both brain and bone marrow were used. For virtual knockout, a prescreening step selected blood-end genes in the top 200 by both mean expression and expression variance (intersection), retaining 45 blood-end candidates while preserving all 74 brain-end hub genes. Bone marrow genotype (5xFAD vs.\ WT) was demultiplexed using UMAP-based clustering on brain-end differentially expressed gene signatures.

\textbf{ADNI data preprocessing.} NULISA plasma proteomics were normalized and log-transformed. Gene expression microarray probes were mapped to gene symbols (grouped by Symbol, averaged across probes). Subject identifiers were converted from PTID to RID using the ADNI REGISTRY file. For survival analysis, the event was defined as MMSE decline $\geq$ 3 points from baseline; subjects without events were censored at the last available visit.

\subsubsection{Neural ODE Framework: Architecture, Training, and Benchmark}

The Neural ODE models a continuous-time dynamical system:
\begin{equation}
\frac{d\mathbf{z}(t)}{dt} = f_\theta\bigl(\mathbf{z}(t),\, t\bigr), \quad \mathbf{z}(0) = \mathbf{x}_{\text{source}}, \quad \mathbf{z}(1) \approx \mathbf{x}_{\text{target}}
\end{equation}
where $f_\theta: \mathbb{R}^d \times [0,1] \to \mathbb{R}^d$ is a 4-layer multilayer perceptron (hidden dimension 256, SiLU activation) parameterizing the vector field, $\mathbf{x}_{\text{source}}$ is a cell from the brain (or blood) distribution, and $\mathbf{z}(1)$ is the predicted state in the target compartment. The ODE is integrated using the RK4 (Runge--Kutta 4th order) solver with 4 sub-steps:
\begin{align}
\mathbf{k}_1 &= f_\theta(\mathbf{z}_n,\, t_n) \\
\mathbf{k}_2 &= f_\theta(\mathbf{z}_n + \tfrac{h}{2}\mathbf{k}_1,\, t_n + \tfrac{h}{2}) \\
\mathbf{k}_3 &= f_\theta(\mathbf{z}_n + \tfrac{h}{2}\mathbf{k}_2,\, t_n + \tfrac{h}{2}) \\
\mathbf{k}_4 &= f_\theta(\mathbf{z}_n + h\mathbf{k}_3,\, t_n + h) \\
\mathbf{z}_{n+1} &= \mathbf{z}_n + \tfrac{h}{6}(\mathbf{k}_1 + 2\mathbf{k}_2 + 2\mathbf{k}_3 + \mathbf{k}_4)
\end{align}
where $h$ is the step size.

\textbf{Sinkhorn OT loss.} The training objective minimizes the entropic-regularized optimal transport distance between the predicted batch $\{\mathbf{z}_i(1)\}_{i=1}^B$ and a target batch $\{\mathbf{x}_j^{\text{target}}\}_{j=1}^B$:
\begin{equation}
\mathcal{L}_{\text{OT}} = \langle \mathbf{P}^*, \mathbf{C} \rangle, \quad \mathbf{P}^* = \arg\min_{\mathbf{P} \in \Pi(\mathbf{a},\mathbf{b})} \langle \mathbf{P}, \mathbf{C} \rangle - \epsilon \, H(\mathbf{P})
\end{equation}
where $\mathbf{C}_{ij} = \|\mathbf{z}_i(1) - \mathbf{x}_j^{\text{target}}\|^2$ is the cost matrix, $\mathbf{a},\mathbf{b}$ are uniform marginals, $H(\mathbf{P}) = -\sum_{ij} P_{ij}(\log P_{ij} - 1)$ is the entropy, $\epsilon$\,=\,0.15 is the regularization strength, and $\Pi(\mathbf{a},\mathbf{b})$ is the transportation polytope. The optimal coupling $\mathbf{P}^*$ is computed via Sinkhorn iterations:
\begin{equation}
\mathbf{u}^{(k+1)} = \frac{\mathbf{a}}{\mathbf{K}\mathbf{v}^{(k)}}, \quad \mathbf{v}^{(k+1)} = \frac{\mathbf{b}}{\mathbf{K}^\top\mathbf{u}^{(k+1)}}, \quad \mathbf{K} = e^{-\mathbf{C}/\epsilon}
\end{equation}
with 50 iterations. The total loss additionally includes an $L_2$ regularization term: $\mathcal{L} = \mathcal{L}_{\text{OT}} + \lambda \|\theta\|_2^2$ ($\lambda$\,=\,$10^{-3}$). Training: 200 transport pairs per epoch, batch size 128, Adam optimizer (lr\,=\,$2 \times 10^{-4}$), 100 epochs for transcriptomics/proteomics.

The benchmark compared four methods on held-out validation data ($n$\,=\,1{,}400 evaluation samples): (1) Neural ODE, (2) Direct OT Map (Sinkhorn barycentric projection), (3) Linear Ridge Regression, (4) Identity baseline. Mapping quality was assessed by OT distance and gene--gene correlation-structure error. World-model-exclusive capabilities were quantified by Jacobian AUC (concordance with GenKI functional validation, $n$\,=\,15 pairs), continuous trajectory monotonicity ($t \in [0,1]$), and round-trip reversibility (Pearson $r$).

\subsubsection{Cross-Species Transfer Validation}
To test whether the mouse-trained world model captures evolutionarily conserved brain--blood regulatory structure, we loaded the trained transcriptomics model and evaluated it on human ADNI bulk gene expression microarray data (745 subjects, 49{,}386 probes collapsed to gene-level by median). ADNI expression was z-scored per gene across subjects and aligned to the 2,000-gene model space (1{,}961 genes matched, 98\%). We measured: (1) the OT distance between model input (human blood expression) and model output (predicted brain-end state), as a non-triviality check that the model transforms rather than copies the input; (2) the OT distance between the model's human prediction and real 5xFAD mouse brain expression, as a cross-species alignment metric; and (3) the effect of light fine-tuning (20 epochs, learning rate $10^{-5}$, $L_2$ regularization toward pretrained weights) on both metrics, with concurrent monitoring of mouse OT distance to detect catastrophic forgetting.

\subsubsection{Jacobian Sensitivity Analysis}

The Jacobian matrix of the trained Neural ODE vector field quantifies the sensitivity of each target-tissue molecule to perturbations in each source-tissue molecule:
\begin{equation}
\mathbf{J}_{ij} = \frac{\partial z_j(1)}{\partial z_i(0)} = \frac{\partial (\text{blood}_j)}{\partial (\text{brain}_i)}
\end{equation}
This was computed via automatic differentiation (PyTorch \texttt{autograd}) of the RK4-integrated trajectory at the source distribution mean. For a 4-step RK4 integrator, the Jacobian accumulates through the chain rule:
\begin{equation}
\frac{\partial \mathbf{z}_{n+1}}{\partial \mathbf{z}_n} = \mathbf{I} + \frac{h}{6}\left(\frac{\partial \mathbf{k}_1}{\partial \mathbf{z}_n} + 2\frac{\partial \mathbf{k}_2}{\partial \mathbf{z}_n} + 2\frac{\partial \mathbf{k}_3}{\partial \mathbf{z}_n} + \frac{\partial \mathbf{k}_4}{\partial \mathbf{z}_n}\right)
\end{equation}
and the full Jacobian is the product across all integration steps: $\mathbf{J} = \prod_{n=0}^{N-1} \frac{\partial \mathbf{z}_{n+1}}{\partial \mathbf{z}_n}$. Three confidence metrics were computed from 100 bootstrap replicates: stability (fraction of bootstraps retaining the edge), signal-to-noise ratio (mean Jacobian / SD across bootstraps), and directional consistency (agreement between forward and reverse Jacobian signs). Edges in the top 5\% by absolute Jacobian magnitude and passing confidence thresholds were retained.

\subsubsection{Hub Gene Identification via MultiXrank}

MultiXrank \citep{baptista2022} was applied to multi-layer networks comprising the Jacobian sensitivity edges (brain and blood compartments per omics layer) and STRING protein interaction networks. The algorithm performs random walk with restart on the multiplex network, yielding per-gene centrality scores. Genes in the top 10th percentile of centrality were designated as hubs (393 transcriptomic, 21 proteomic brain-end hubs).

\subsubsection{Virtual Knockout}

\textbf{GenKI.} For each omics layer, a PCNet gene regulatory network was constructed from the single-cell expression matrix using the method of \citep{senuma2023}, with quantile normalization and $q$\,=\,0.85 (proteomics) or 0.85 (transcriptomics) edge retention. A variational graph autoencoder (VGAE) was trained on the wild-type network to learn a latent representation. Virtual knockout was performed by zeroing the target gene's expression, re-training the VGAE, and computing the Kullback--Leibler divergence between wild-type and knockout latent distributions. Significance was assessed via comparison to expression-matched negative control genes (Cohen's $d$, permutation $p$-value). The transcriptomics VGAE used a transcription-group-specific parameter set (learning rate $10^{-4}$, out\_channels\,=\,16, $\beta$\,=\,0.1, 200 metacells) calibrated for the larger gene set.

\textbf{Tissue-specific PPI propagation.} Brain- and blood-specific PPI networks were loaded from the tissue-specific PPI atlas \citep{lamantrip2025} (12{,}002 brain, 10{,}149 blood proteins; association score threshold 0.5). Random walk with restart ($\alpha$\,=\,0.3) was performed on the tissue-specific network before and after node removal (graph mutilation). The effect score for each observation protein was the relative change in steady-state probability between intact and knockout networks.

\textbf{SCENIC GRN perturbation.} A gene regulatory network (GRN) was inferred from the 5xFAD single-cell expression matrix (3{,}000 cells $\times$ 2{,}078 genes) using GRNBoost2 \citep{moerman2019}, which fits a gradient-boosted regression model ($n$\,=\,50 estimators, max depth 2) predicting each target gene from all other genes. Regulatory edges with feature importance $>$\,0.005 were retained. Regulons were constructed by grouping each transcription factor's top-50 importance targets (requiring $\geq$5 targets per regulon), yielding 136 regulons (median size 8, range 5--38). Baseline regulon activity was scored using AUCell \citep{aibar2017}, which computes the area under the recovery curve of regulon genes ranked by expression in each cell. Virtual knockout was performed by zeroing the target gene's expression, recomputing AUCell scores, and measuring the mean absolute change in regulon activity ($\Delta$) across all 136 regulons. A $Z$-score was computed as $Z = (\Delta_i - \bar{\Delta}) / \sigma_\Delta$ for each regulon $i$; $|Z| > 1.96$ defined significant perturbation. The same set of KO genes nominated by GenKI was used, providing methodological cross-validation through an independent perturbation principle (GRN regulon activity shift vs.\ VGAE latent-space displacement).

\subsubsection{Diagnostic Panel Construction}

VK-nominated blood-end biomarker candidates were evaluated using L1-regularized logistic regression (solver: liblinear, $C$\,=\,0.1) with 5-fold stratified cross-validation repeated 5 times. Area under the ROC curve (AUC) was the primary metric. Single-gene AUC was computed for all candidates.

\subsubsection{Pathway Enrichment Analysis of VK Perturbation Targets}

To translate virtual-knockout perturbation sets into biological-process interpretation, we performed Gene Ontology Biological Process over-representation analysis (GO-BP ORA) using \texttt{gseapy} (v1.3.0, Enrichr backend \citep{kuleshov2016}) on the union of significantly perturbed targets across all knockouts per method and direction. For each VK method (PPI propagation, GenKI proteomics, GenKI transcriptomics) and direction (forward/reverse), the perturbation target set was defined as all genes with $Z_\text{score} > 2$ in any KO's per-target ranking (\texttt{*\_gene\_ranking.csv}). ORA was run against the \texttt{GO\_Biological\_Process\_2023} gene-set library (Enrichr); the top-12 terms by Benjamini--Hochberg-adjusted $p$-value were retained per analysis. SCENIC regulon rankings were excluded because their per-regulon (not per-gene) score schema is incompatible with gene-level ORA. Six method$\times$direction combinations were analyzed in total.

For the Sankey--bubble visualization (Figure~\ref{fig:vk}i), the gene$\rightarrow$pathway edgelist was constructed by intersecting each enriched pathway's gene list with the PPI perturbation target set of the corresponding direction, keeping up to 8 genes per pathway for readability. The Sankey flow (gene$\rightarrow$pathway) was rendered with \texttt{ggsankeyfier} (v0.1.8); the aligned bubble panel encoded $-\log_{10}(\text{adjusted}\ p)$ on the $x$-axis, gene-count as bubble size, and overlap hit-ratio as color, with each pathway's $y$-position matched to its Sankey node coordinate.

\subsubsection{Aging Pseudotime Analysis}

To characterize the dynamics of validated target expression across the lifespan, we computed a continuous aging pseudotime from the Tabula Muris Senis (TMS) brain single-cell RNA-seq dataset \citep{schaum2019} (5{,}000 brain cells across ages 3m/18m/24m) using the scanpy (v1.11.5) diffusion map and diffusion pseudotime (DPT) pipeline \citep{wolf2018,haghverdi2016}. Cells were filtered (genes expressed in $\geq$10 cells), normalized to 10{,}000 counts, log-transformed, and the top 2{,}500 highly variable genes (seurat\_v3) were selected. PCA (30 components) fed a $k$-nearest-neighbor graph ($k=20$), from which a diffusion map (15 components) was computed; DPT was rooted at the 3-month-cell centroid in diffusion space. Because the DPT root can land at either end of the diffusion axis, the sign of pseudotime was flipped if its Spearman correlation with chronological age was negative, so that 3m = young (low) and 24m = old (high) by construction.

Human validated targets were mapped to mouse one-to-one orthologs using the MGI HOM\_MouseHumanSequence.rpt database (17{,}609 pairs); 26 of 31 candidate targets had a measurable ortholog in the TMS brain data (5 dropped: CYP4F3, HBA1, MCTP2, TAFA5, VEGFD, no 1:1 ortholog or not expressed). Ortholog expression (log-normalized) was projected onto pseudotime, $z$-scored across cells, binned into 30 pseudotime quantiles, and the per-bin mean was rendered as a 3D surface plot ($x$ = pseudotime, $y$ = gene, $z$ = mean $z$-expression; Figure~\ref{fig:diagnostic}h).

\subsubsection{Survival Analysis}

Cox proportional hazards regression was performed using the \texttt{lifelines} library. The event was defined as MMSE decline $\geq$ 3 points from baseline screening visit. Follow-up time was calculated from baseline to event or last visit. All models adjusted for APOE4 carrier status (binary). Gene expression was Z-scored. For diagnostic biomarkers (forward direction), survival significance indicates prognostic value (prediction of cognitive trajectory). For therapeutic targets (reverse direction), survival significance indicates that the target's expression level influences disease progression (therapeutic response prediction).

\subsubsection{Drug Target Mining}

Multi-database drug mining was performed using OpenTargets (known drugs, clinical phase, tractability), ChEMBL (mechanism of action), DGIdb (literature-mined interactions), and PubChem (bioactivity screening data). The composite Drug Evidence Score weighted clinical phase (45\%), tractability (20\%), DGIdb evidence (10\%), ChEMBL evidence (10\%), and PubChem activity (15\%). Targets were tiered as Approved ($\geq$1 approved drug), Late Clinical (Phase III), Early Clinical (Phase II), or Tractable.

\subsection*{Quantification and Statistical Analysis}

All statistical analyses were performed in Python (version 3.10). Cross-tissue mapping quality was assessed by Sinkhorn OT distance and correlation-structure error. Jacobian concordance was measured by AUC. Virtual knockout significance used Cohen's $d$ and permutation testing against expression-matched negative controls. Diagnostic performance was assessed by cross-validated AUC. Survival analysis used Cox proportional hazards regression with APOE4 adjustment. All $p$-values are two-sided. Multiple testing correction was not applied to the discovery-phase VK analysis (exploratory); clinical validation used $p < 0.05$ as the significance threshold.

% ======================================================================
\begin{thebibliography}{99}

\bibitem{chen2018} Chen, R.T.Q. et al. (2018). Neural ordinary differential equations. \textit{NeurIPS}.

\bibitem{cuturi2013} Cuturi, M. (2013). Sinkhorn distances: lightspeed computation of optimal transport. \textit{NeurIPS}.

\bibitem{ha2018} Ha, D. \& Schmidhuber, J. (2018). World models. \textit{arXiv:1803.10122}.

\bibitem{oakley2006} Oakley, H. et al. (2006). Intraneuronal $\beta$-amyloid aggregates, neurodegeneration, and neuron loss in transgenic mice with five familial Alzheimer's disease mutations: potential factors in amyloid plaque formation. \textit{J.\ Neurosci.} 26, 10129--10140.

\bibitem{jawhar2012} Jawhar, S. et al. (2012). Motor deficits, neuron loss, and reduced anxiety coinciding with axonal degeneration and intraneuronal A$\beta$ aggregation in the 5XFAD mouse model of Alzheimer's disease. \textit{Neurobiol.\ Aging} 33, 196.e29--196.e40.

\bibitem{eppig2015} Eppig, J.T. et al. (2015). The Mouse Genome Database (MGD): facilitating mouse as a model for human biology and disease. \textit{Nucleic Acids Res.} 43, D726--D736.

\bibitem{aibar2017} Aibar, S. et al. (2017). SCENIC: single-cell regulatory network inference and clustering. \textit{Nat.\ Methods} 14, 1083--1086.

\bibitem{moerman2019} Moerman, T. et al. (2019). GRNBoost2 and Arboreto: efficient and scalable inference of gene regulatory networks. \textit{Bioinformatics} 35, 2159--2161.

\bibitem{lamantrip2025} Laman Trip, D.S. et al. (2025). A tissue-specific atlas of protein--protein associations enables prioritization of candidate disease genes. \textit{Nat.\ Biotechnol.} 44, 654--667.

\bibitem{senuma2023} Yang, Y. et al. (2023). Gene knockout inference with variational graph autoencoder learning. \textit{Nucleic Acids Res.} 51, 6578--6592.

\bibitem{baptista2022} Baptista, A. et al. (2022). Universal multilayer network exploration by random walk with restart. \textit{Commun.\ Phys.} 5, 170.

\bibitem{bettcher2021} Bettcher, B.M. et al. (2021). Peripheral and central immune system crosstalk in Alzheimer disease --- a research prospectus. \textit{Nat.\ Rev.\ Neurol.} 17, 689--701.

\bibitem{gate2024} Gate, D. et al. (2020). Clonally expanded CD8 T cells patrol the cerebrospinal fluid in Alzheimer's disease. \textit{Nature} 577, 399--404.

\bibitem{xu2023} Xu, H. et al. (2021). Single-cell RNA sequencing of peripheral blood reveals immune cell signatures in Alzheimer's disease. \textit{Front.\ Immunol.} 12, 645666.

\bibitem{sweeney2018} Sweeney, M.D. et al. (2018). The role of brain vasculature in neurodegenerative disorders. \textit{Nat.\ Neurosci.} 21, 1318--1331.

\bibitem{nation2019} Nation, D.A. et al. (2019). Blood--brain barrier breakdown is an early biomarker of human cognitive dysfunction. \textit{Nat.\ Med.} 25, 270--276.

\bibitem{venegas2017} Venegas, C. et al. (2017). Microglia-derived ASC specks cross-seed amyloid-$\beta$ in Alzheimer's disease. \textit{Nature} 552, 355--361.

\bibitem{hessel2023} Heneka, M.T., McManus, R.M. \& Latz, E. (2018). Inflammasome signalling in brain function and neurodegenerative disease. \textit{Nat.\ Rev.\ Neurosci.} 19, 610--621.

\bibitem{wang2022} Kreitmaier, P. et al. (2023). Insights from multi-omics integration in complex disease primary tissues. \textit{Trends Genet.} 39, 46--58.

\bibitem{heppner2015} Heppner, F.L., Ransohoff, R.M. \& Becher, B. (2015). Immune attack: the role of inflammation in Alzheimer disease. \textit{Nat.\ Rev.\ Neurosci.} 16, 358--372.

\bibitem{viggars2021} Bell, R.D. \& Zlokovic, B.V. (2009). Neurovascular mechanisms and blood--brain barrier disorder in Alzheimer's disease. \textit{Acta Neuropathol.} 118, 103--113.

\bibitem{boza2017} Boza-Serrano, A. et al. (2019). Galectin-3, a novel endogenous TREM2 ligand, detrimentally regulates inflammatory response in Alzheimer's disease. \textit{Acta Neuropathol.} 138, 251--273.

\bibitem{bunne2024} Bunne, C. et al. (2024). How to build the virtual cell with artificial intelligence: priorities and opportunities. \textit{Cell} 187, 7045--7063.

\bibitem{lorenzl2003} Lorenzl, S. et al. (2003). Increased plasma levels of matrix metalloproteinase-9 in patients with Alzheimer's disease. \textit{Neurochem.\ Int.} 43, 191--196.

\bibitem{yang2007} Yang, Y., Estrada, E.Y., Thompson, J.F., Liu, W. \& Rosenberg, G.A. (2007). Matrix metalloproteinase-mediated disruption of tight junction proteins in cerebral vessels is reversed by synthetic matrix metalloproteinase inhibitor in focal ischemia in rat. \textit{J.\ Cereb.\ Blood Flow Metab.} 27, 697--709.

\bibitem{tsai2007} Tsai, K.J. et al. (2007). G-CSF rescues the memory impairment of animal models of Alzheimer's disease. \textit{J.\ Exp.\ Med.} 204, 1273--1280.

\bibitem{zenaro2015} Zenaro, E. et al. (2015). Neutrophils promote Alzheimer's disease-like pathology and cognitive decline via LFA-1 integrin. \textit{Nat.\ Med.} 21, 880--886.

\bibitem{dudvarzi2025} Camprub\'i-Ferrer, L. et al. (2025). Galectin-3 deletion modulates microglial phenotype and A$\beta$ response via TREM2 activation while attenuating neuroinflammation. \textit{bioRxiv} 2025.03.17.643790.

\bibitem{bakshi2011} Bakshi, P. et al. (2011). Depletion of CXCR2 inhibits $\gamma$-secretase activity and amyloid-$\beta$ production in a murine model of Alzheimer's disease. \textit{Cytokine} 53, 163--169.

\bibitem{ryu2015} Ryu, J.K. et al. (2015). Pharmacological antagonism of interleukin-8 receptor CXCR2 inhibits inflammatory reactivity and is neuroprotective in an animal model of Alzheimer's disease. \textit{J.\ Neuroinflammation} 12, 144.

\bibitem{trem1review2024} Wilson, E.N., Wang, C., Swarovski, M.S. et al. (2024). TREM1 disrupts myeloid bioenergetics and cognitive function in aging and Alzheimer disease mouse models. \textit{Nat.\ Neurosci.} 27, 873--885.

\bibitem{hansson2024} Palmqvist, S., Tideman, P., Mattsson-Carlgren, N. et al. (2025). Plasma phospho-tau217 for Alzheimer's disease diagnosis in primary and secondary care using a fully automated platform. \textit{Nat.\ Med.} 31, 2036--2043.

\bibitem{barthelemy2025} Horie, K., Salvad\'o, G., Barth\'elemy, N.R. et al. (2025). Plasma MTBR-tau243 biomarker identifies tau tangle pathology in Alzheimer's disease. \textit{Nat.\ Med.} 31, 2044--2053.

\bibitem{mmp9mouse2021} Ringland, C. et al. (2021). MMP9 modulation improves specific neurobehavioral deficits in a mouse model of Alzheimer's disease. \textit{BMC Neurosci.} 22, 39.

\bibitem{li2025mmp} Li, S. (2025). Matrix metalloproteinases as therapeutic targets of neuroinflammation in Alzheimer's disease. \textit{Biomed.\ J.\ Sci.\ Tech.\ Res.} 63, 38014--38017.

\bibitem{kuleshov2016} Kuleshov, M.V.\ et al.\ (2016). Enrichr: a comprehensive gene set enrichment analysis web server. \textit{Nucleic Acids Res.} 44, W90--W97.

\bibitem{schaum2019} Schaum, N.\ et al.\ (2020). Ageing hallmarks exhibit organ-specific temporal signatures. The Tabula Muris Senis Consortium. \textit{Nature} 583, 596--602.

\bibitem{wolf2018} Wolf, F.A.\ et al.\ (2018). SCANPY: large-scale single-cell gene expression data analysis. \textit{Genome Biol.} 19, 15.

\bibitem{haghverdi2016} Haghverdi, L., Buettner, M., Wolf, F.A. \& Theis, F.J. (2016). Diffusion pseudotime robustly reconstructs lineage branching. \textit{Nat.\ Methods} 13, 845--848.

\end{thebibliography}

% ======================================================================
\section*{Declarations}

\section*{Competing Interests}
The authors declare no competing interests.

\section*{Author Contributions}
Z.D.\ designed and implemented the computational framework, performed all analyses, and wrote the manuscript. X.M.\ contributed to the world-model conceptualization and supervised the study. L.W.\ supervised the overall research direction and secured funding. All authors edited and approved the final manuscript.

\section*{Ethics statement}
This study is a secondary computational analysis of previously published, publicly available, and de-identified datasets. The 5xFAD mouse single-cell RNA-seq data (GSE329430) were generated under protocols approved by the original data generators' institutional animal ethics committees. ADNI data were collected in accordance with the ethical standards of the institutional review boards of all participating ADNI sites after obtaining written informed consent from all participants or their authorised representatives; all ADNI data accessed here are de-identified. No new experiments involving human participants, human tissue, or live animals were conducted for this study.

\section*{AI-use disclosure}
In preparing this manuscript, the authors used large language models (Z.ai GLM and OpenAI GPT) as writing assistants for code drafting, language editing, and formatting of the manuscript text and figures. All scientific content, analyses, interpretations, and conclusions were conceived, verified, and approved by the authors, who take full responsibility for the integrity and accuracy of all reported results.

\section*{Acknowledgments}
Data collection and sharing for this project was funded by the Alzheimer's Disease Neuroimaging Initiative (ADNI) (National Institutes of Health Grant U01 AG024904). We thank the Tabula Muris Senis Consortium and the GEO data contributors for making single-cell datasets publicly available.

\section*{Abbreviations}
AD: Alzheimer's disease; APOE4: apolipoprotein E4; AUC: area under the curve; CV: cross-validation; DPT: diffusion pseudotime; GO-BP: Gene Ontology Biological Process; GRN: gene regulatory network; HVG: highly variable gene; KL: Kullback--Leibler; MGI: Mouse Genome Informatics; MMP9: matrix metalloproteinase-9; Neural ODE: neural ordinary differential equation; ORA: over-representation analysis; OT: optimal transport; PPI: protein--protein interaction; SCENIC: single-cell regulatory network inference and clustering; TMS: Tabula Muris Senis; VK: virtual knockout; VGAE: variational graph autoencoder; 5xFAD: five familial AD mutations transgenic mouse.

\section*{Keywords}
Alzheimer's disease; world model; Neural ODE; virtual knockout; cross-tissue mapping; drug repurposing; diagnostic biomarker; Jacobian sensitivity; pathway enrichment; aging pseudotime.

\end{document}
